% Created 2019-10-14 Пн 16:47
% Intended LaTeX compiler: pdflatex
\documentclass[11pt]{article}
\usepackage[utf8]{inputenc}
\usepackage[T1]{fontenc}
\usepackage{graphicx}
\usepackage{grffile}
\usepackage{longtable}
\usepackage{wrapfig}
\usepackage{rotating}
\usepackage[normalem]{ulem}
\usepackage{amsmath}
\usepackage{textcomp}
\usepackage{amssymb}
\usepackage{capt-of}
\usepackage{hyperref}
\usepackage[T2A]{fontenc}
\usepackage[a4paper,left=3cm,top=2cm,right=1.5cm,bottom=2cm,marginparsep=7pt,marginparwidth=.6in]{geometry}
\usepackage{cmap}
\usepackage{xcolor}
\usepackage{listings}
\usepackage{polyglossia}
\setdefaultlanguage{russian} \setotherlanguage{english}
\setmainfont{Liberation Serif}
\setsansfont{Liberation Sans}
\setmonofont[Contextuals=Alternate,Ligatures={TeX}]{Fira Code Regular}
\author{Krutko Nikita / KrutNA}
\date{\today}
\title{}
\hypersetup{
 pdfauthor={Krutko Nikita / KrutNA},
 pdftitle={},
 pdfkeywords={},
 pdfsubject={},
 pdfcreator={Emacs 26.1 (Org mode 9.1.9)}, 
 pdflang={Russian}}
\begin{document}

\large
\thispagestyle{empty}
\begin{center}
\textbf{Национальный Исследовательский Университет ИТМО}\\
\textbf{Факультет Программной Инженерии и Компьютерной Техники}\\
\end{center}
\vspace{2em}
\begin{center}
\includegraphics[width=120pt]{../itmo-logo.png}
\end{center}
\LARGE
\vspace{5em}
\begin{center}
\textbf{Вариант № 14 \% 9 = 5}\\
\textbf{Лабораторная работа № 3}\\
\Large
\textbf{по дисциплине}\\
\LARGE
\textbf{\emph{'Информатика'}}\\
\end{center}
\vspace{11em}
\large
\begin{flushright}
\textbf{Выполнил:}\\
\textbf{Студент группы P3113}\\
\textbf{\emph{Крутько Никита;} : 242570}\\
\textbf{Преподаватель:}\\
\textbf{\emph{Малышева Татьяна Алексеевна}}\\
\end{flushright}
\vspace{4em}
\large
\begin{center}
\textbf{Санкт-Петербург 2019 г.}
\end{center}
\pagebreak{}
\setcounter{tocdepth}{2}
\tableofcontents
\normalsize
\section{Задание}
\label{sec:org8e82179}
\subsection{Файл}
\label{sec:orgdd4dc55}
\begin{itemize}
\item Создать следующего вида исходный файл из десяти строк, содержащий в каждой строке ФИО, дату рождения, дату получения паспорта и баллы ЕГЭ по трём предметам:
\end{itemize}
\scriptsize
\begin{verbatim}
Апельсинов А.А. 08.02.2000 17.03.2014 90 100 91
Виноградов В.В. 09.03.1999 15.04.2013 67 99 98
Яблоков Я.Я. 10.04.2000 19.05.2014 94 94 94
Морковкин М.М. 11.05.1999 17.06.2013 91 82 73
Жиноградов В.В. 09.03.1999 15.04.2013 67 99 98
Сяблоков Я.Я. 10.04.2000 19.05.2014 94 94 94
Рорковкин М.М. 11.05.1999 17.06.2013 91 82 73
Диноградов В.В. 09.03.1999 15.04.2013 67 99 98
Няблоков Я.Я. 10.04.2000 19.05.2014 94 94 94
Корковкин М.М. 11.05.1999 17.06.2013 91 82 73
\end{verbatim}

\subsection{Заданиеи}
\label{sec:orgf128ebb}
\normalsize
\begin{itemize}
\item Не используя готовые сторонние подключаемые функции для факториала, int(), bin() и т.п., написать программу на языке Python 3.x, которая бы вычисляла среднее значение баллов ЕГЭ, сортировала строки по указанной колонке в обратном порядке (от большего к меньшему) и выводила результат следующего вида (для сортировки по дате рождения):
\end{itemize}
\scriptsize
\begin{verbatim}
Яблоков Я.Я. | 10.04.2000 | 19.05.2014 | 94 94 94 -> 94.000000
Сяблоков Я.Я. | 10.04.2000 | 19.05.2014 | 94 94 94 -> 94.000000
Рорковкин М.М. | 11.05.1999 | 17.06.2013 | 91 82 73 -> 82.000000
Няблоков Я.Я. | 10.04.2000 | 19.05.2014 | 94 94 94 -> 94.000000
Морковкин М.М. | 11.05.1999 | 17.06.2013 | 91 82 73 -> 82.000000
Корковкин М.М. | 11.05.1999 | 17.06.2013 | 91 82 73 -> 82.000000
Жиноградов В.В. | 09.03.1999 | 15.04.2013 | 67 99 98 -> 88.000000
Диноградов В.В. | 09.03.1999 | 15.04.2013 | 67 99 98 -> 88.000000
Виноградов В.В. | 09.03.1999 | 15.04.2013 | 67 99 98 -> 88.000000
Апельсинов А.А. | 08.02.2000 | 17.03.2014 | 90 100 91 -> 93.666667
\end{verbatim}

\section{Код Python}
\label{sec:org72fd0fa}
\scriptsize
\lstset{language=Python,label= ,caption= ,captionpos=b,numbers=none}
\begin{lstlisting}
#!/usr/bin/env python3
def get_avg(array):
    return sum(int(i) for i in array)/len(array)


def split(line):
    res = line.split(' ')
    return [' '.join(res[0:2]), res[2], res[3],
	    ' '.join(res[4:]), get_avg(res[4:])]


def sort(array, col=4, ascending=True):
    return sorted(array, key=lambda val: val[col], reverse=not ascending)


def print_formated(line):
    print("%s -> %f" % (' | '.join(line[:4]), line[4]))


if __name__ == '__main__':
    file = "file"
    with open(file, "r") as f:
	result = list(map(lambda line: split(line.rstrip()), f.readlines()))
    result = sort(result, 0, False)
    list(map(print_formated, result))
\end{lstlisting}

\section{Вывод}
\label{sec:orgafd301a}
\normalsize
Немножко кода на питоне, ничего особенного
\end{document}